\documentclass[11pt,a4paper]{article}
\usepackage[utf8]{inputenc}
\usepackage[T1]{fontenc}
\usepackage[french]{babel}
\usepackage{geometry}
\usepackage{hyperref}
\usepackage{enumitem}
\usepackage{listings}
\geometry{margin=2.2cm}
\hypersetup{colorlinks=true,linkcolor=blue,urlcolor=blue}

\lstdefinestyle{cmd}{
  basicstyle=\ttfamily\small,
  breaklines=true,
  frame=single,
  columns=fullflexible,
  keepspaces=true
}

\title{Rapport d'equipe\\Ajout d'une nouvelle serie STM32Cube dans le pipeline}
\author{Projet STM32Cube Preprocessing}
\date{\today}

\begin{document}
\maketitle
\tableofcontents
\newpage

\section{Objet du rapport}
Ce document formalise la procedure standard d'ajout d'une nouvelle serie STM32Cube
(exemple: C5) dans le pipeline. Il precise les fichiers a modifier, la signification
des champs de configuration, la source des donnees a collecter, l'integration dans
l'automatisation, ainsi que les etapes de validation avant mise en production.

\section{Perimetre technique}
Le perimetre couvre:
\begin{itemize}[leftmargin=*]
  \item la configuration serie,
  \item l'execution du pipeline (Prepare/Upload),
  \item le split des issues par composant,
  \item l'upload vers la KB,
  \item l'integration dans les workflows GitHub Actions,
  \item le test sur une autre machine avec self-hosted runner.
\end{itemize}

\section{Fichiers a modifier}
\subsection{Configuration}
\begin{itemize}[leftmargin=*]
  \item \texttt{shared/config/config\_c5.json} (nouveau fichier)
  \item \texttt{shared/config/config\_all\_series.json} (mise a jour)
\end{itemize}

\subsection{Scripts workflow}
\begin{itemize}[leftmargin=*]
  \item \texttt{pipeline\_Automation/workflow/Run\_Single\_Series\_Full\_Pipeline\_And\_Upload.ps1}
  \item \texttt{pipeline\_Automation/workflow/Split\_Issues\_By\_Component.ps1}
  \item \texttt{pipeline\_Automation/workflow/Upload\_Issues\_By\_Component.ps1}
  \item \texttt{pipeline\_Automation/workflow/run\_all\_series.py}
\end{itemize}

\subsection{Workflows CI/CD}
\begin{itemize}[leftmargin=*]
  \item \texttt{.github/workflows/auto\_update\_kb.yml}
  \item \texttt{.github/workflows/manual\_auto\_update\_kb.yml}
\end{itemize}

\section{Definition des champs de configuration}
\subsection{Fichier \texttt{config\_<serie>.json}}
\textbf{owner}\\
Organisation source des depots (STMicroelectronics).

\textbf{repos}\\
Liste ordonnee des depots de la serie. Le premier element doit etre le repo parent
(exemple: STM32CubeC5), suivi des depots lies (HAL/CMSIS/BSP/Middleware).

\textbf{data\_dir}\\
Dossier des artefacts data intermediaires (valeur standard: \texttt{data}).

\textbf{repos\_dir}\\
Dossier local des clones (valeur standard: \texttt{GitHub\_repos}).

\textbf{auto\_clone\_missing\_repos}\\
Activation du clonage automatique des depots manquants.

\textbf{repo\_clone\_urls}\\
Dictionnaire \texttt{repo -> URL git} pour les depots relies.

\textbf{linked\_repo\_refs}\\
Dictionnaire \texttt{repo -> SHA} correspondant aux refs des sous-modules pour la
release cible. Ce champ est essentiel pour la reproductibilite.

\textbf{commit\_pr\_scope}\\
Strategie de liaison des commits/PR (valeur recommandee: \texttt{linked\_or\_merged}).

\textbf{github\_token\_env}\\
Nom de variable d'environnement contenant le token GitHub.

\textbf{rate\_limits}\\
Temporisations API (issues/comments) pour stabiliser les appels.

\textbf{ingest\_version}\\
Version logique de la configuration d'ingestion.

\subsection{Fichier \texttt{config\_all\_series.json}}
\textbf{all\_series}\\
Liste des series actives.

\textbf{series\_config\_files}\\
Mapping \texttt{Serie -> fichier config}.

\textbf{kb\_datasource\_ids}\\
IDs des datasources par serie (issues, files, diagnostic, resolver).

\textbf{kb\_id}\\
Identifiant de la Knowledge Base cible.

\section{Source des informations: repos, urls, refs}
Les informations sont extraites du repo parent de la serie, positionne sur la
release cible.

\subsection{Commandes de collecte}
\begin{lstlisting}[style=cmd]
git clone --recursive https://github.com/STMicroelectronics/STM32CubeC5.git
cd STM32CubeC5
git checkout <tag_cible>
git submodule update --init --recursive

# URLs des sous-modules
git config -f .gitmodules --get-regexp "submodule\..*\.url"

# SHAs de references
git submodule status --recursive

# Paths des sous-modules
git config -f .gitmodules --get-regexp "submodule\..*\.path"
\end{lstlisting}

\subsection{Regles de remplissage}
\begin{itemize}[leftmargin=*]
  \item \texttt{repos[0]} = repo parent de la serie.
  \item \texttt{repo\_clone\_urls} renseigne les depots a cloner.
  \item \texttt{linked\_repo\_refs} renseigne les SHAs de la release cible.
  \item Exclure les depots non relies au perimetre de la serie.
\end{itemize}

\section{Procedure d'integration d'une nouvelle serie}
\subsection{Etape 1: creation du fichier config serie}
Creer \texttt{shared/config/config\_c5.json} a partir d'un modele proche
(\texttt{config\_c0.json} ou \texttt{config\_g4.json}) et renseigner les champs
\texttt{repos}, \texttt{repo\_clone\_urls}, \texttt{linked\_repo\_refs}.

\subsection{Etape 2: declaration globale}
Mettre a jour \texttt{shared/config/config\_all\_series.json}:
\begin{itemize}[leftmargin=*]
  \item ajout dans \texttt{all\_series},
  \item mapping dans \texttt{series\_config\_files},
  \item entree dans \texttt{kb\_datasource\_ids}.
\end{itemize}

\subsection{Etape 3: mapping scripts}
Ajouter la serie dans tous les \texttt{ValidateSet} et \texttt{seriesMap}:
\begin{itemize}[leftmargin=*]
  \item code serie: \texttt{C5}
  \item slug serie: \texttt{stm32cubec5}
  \item repo parent: \texttt{STM32CubeC5}
  \item config: \texttt{shared/config/config\_c5.json}
\end{itemize}

\subsection{Etape 4: mapping batch Python}
Dans \texttt{run\_all\_series.py}, ajouter \texttt{c5} dans \texttt{SERIES\_CONFIGS}.

\subsection{Etape 5: mapping GitHub Actions}
Ajouter la serie C5 dans:
\begin{itemize}[leftmargin=*]
  \item les listes de series autorisees,
  \item les listes par defaut (si desired),
  \item les normalisations d'inputs.
\end{itemize}

\section{Validation fonctionnelle}
\subsection{Prepare}
\begin{lstlisting}[style=cmd]
powershell -ExecutionPolicy Bypass -File .\pipeline_Automation\workflow\Run_Single_Series_Full_Pipeline_And_Upload.ps1 -Series C5 -Mode Prepare
\end{lstlisting}

\subsection{Split composant}
\begin{lstlisting}[style=cmd]
powershell -ExecutionPolicy Bypass -File .\pipeline_Automation\workflow\Split_Issues_By_Component.ps1 -Series C5
\end{lstlisting}

\subsection{Controle couverture et doublons}
\begin{lstlisting}[style=cmd]
python -m pipeline.delivery.validate_issues_by_component_split --summary "datasets/07_delivery/st_ready/by_series/stm32cubec5/issues_json/by_component/summary_by_component_c5.json"
\end{lstlisting}

\subsection{Upload composant}
\begin{lstlisting}[style=cmd]
powershell -ExecutionPolicy Bypass -File .\pipeline_Automation\workflow\Upload_Issues_By_Component.ps1 -Series C5 -DryRun -RemoteUser "<email_st>"
powershell -ExecutionPolicy Bypass -File .\pipeline_Automation\workflow\Upload_Issues_By_Component.ps1 -Series C5 -RemoteUser "<email_st>"
\end{lstlisting}

\section{Test sur une autre machine (self-hosted runner)}
\subsection{Prerequis}
\begin{itemize}[leftmargin=*]
  \item Windows 10/11
  \item Git
  \item Python 3.12+
  \item acces reseau GitHub/ST
\end{itemize}

\subsection{Installation projet}
\begin{lstlisting}[style=cmd]
git clone https://github.com/<org-ou-user>/PFE_Chatbot_STM32Cube.git
cd PFE_Chatbot_STM32Cube
python -m venv .venv
.\.venv\Scripts\Activate.ps1
python -m pip install --upgrade pip
pip install -r requirements.txt
\end{lstlisting}

\subsection{Runner et configuration}
\begin{itemize}[leftmargin=*]
  \item Declarer un self-hosted runner depuis les settings du repo.
  \item Lancer le runner en service.
  \item Verifier le label \texttt{self-hosted}.
  \item Activer \texttt{git core.longpaths=true}.
  \item Configurer secrets: \texttt{ST\_CHATGPT\_API\_KEY}, \texttt{ST\_REMOTE\_USER}.
\end{itemize}

\subsection{Test de bout en bout}
\begin{lstlisting}[style=cmd]
powershell -ExecutionPolicy Bypass -File .\pipeline_Automation\workflow\Run_Single_Series_Full_Pipeline_And_Upload.ps1 -Series C5 -Mode Full -EnableIssuesByComponentFlow -ComponentIssuesOnly -RemoteUser "<email_st>"
\end{lstlisting}

\section{Mode migration vers datasources par composant}
Pour remplacer le datasource global issues par des datasources
\texttt{issues\_<SERIE>\_<COMPONENT>}:
\begin{itemize}[leftmargin=*]
  \item activer \texttt{EnableIssuesByComponentFlow},
  \item activer \texttt{ComponentIssuesOnly}.
\end{itemize}

\section{Checklist Go/No-Go}
\begin{itemize}[leftmargin=*]
  \item config serie creee et controlee,
  \item config globale mise a jour,
  \item mappings scripts a jour,
  \item mappings workflows a jour,
  \item phase Prepare validee,
  \item split composant valide,
  \item validation anti-doublons OK,
  \item upload dry-run OK,
  \item upload reel OK,
  \item datasource IDs traces dans \texttt{config\_all\_series.json}.
\end{itemize}

\section{Points d'attention}
\begin{itemize}[leftmargin=*]
  \item SHAs incorrects dans \texttt{linked\_repo\_refs} => non reproductibilite.
  \item serie declaree en config mais absente des \texttt{ValidateSet/seriesMap} => echec d'execution.
  \item serie absente des workflows => impossible a lancer via dispatch.
  \item upload composant sans split prealable => summary absent.
\end{itemize}

\end{document}
